\documentclass[10pt,a4paper]{amsart}
\usepackage{amsmath,amssymb,amsthm,graphicx}
% Retained preamble for the Stacks Project chapter {CH} (shared project preamble).
% The chapter file pinned-original-source.tex only carries \input{{preamble}} on its first line,
% so no preamble text lives inside the chapter itself.  The Stacks Project uses one single shared
% preamble for every chapter: the snapshot stages it three times (sources/stacks-chow/preamble.tex,
% sources/stacks-etale-cohomology/preamble.tex and focused-seeds/AGP/sources/stacks-snapshot-20261011/
% preamble.tex) and those copies are identical in size.  The block below is copied mechanically,
% line for line, from sources/stacks-chow/preamble.tex lines 14-18 and 154-217: the Xy-pic setup and
% the native theorem declarations together with the mathematical macros.  Nothing was added, removed
% or redefined; the wrapper class and the amsmath/amssymb/amsthm/graphicx base packages are supplied
% by the assembler.
% For commutative diagrams we use Xy-pic
% Chapter: stacks-cohomology.tex -- Cohomology of Sheaves.
\usepackage[all]{xy}

% We use 2cell for 2-commutative diagrams.
\xyoption{2cell}

% Theorem environments.
%
\theoremstyle{plain}
\newtheorem{theorem}[subsection]{Theorem}
\newtheorem{proposition}[subsection]{Proposition}
\newtheorem{lemma}[subsection]{Lemma}

\theoremstyle{definition}
\newtheorem{definition}[subsection]{Definition}
\newtheorem{example}[subsection]{Example}
\newtheorem{exercise}[subsection]{Exercise}
\newtheorem{situation}[subsection]{Situation}

\theoremstyle{remark}
\newtheorem{remark}[subsection]{Remark}
\newtheorem{remarks}[subsection]{Remarks}

\numberwithin{equation}{subsection}

% Macros
%
\def\lim{\mathop{\mathrm{lim}}\nolimits}
\def\colim{\mathop{\mathrm{colim}}\nolimits}
\def\Spec{\mathop{\mathrm{Spec}}}
\def\Hom{\mathop{\mathrm{Hom}}\nolimits}
\def\Ext{\mathop{\mathrm{Ext}}\nolimits}
\def\SheafHom{\mathop{\mathcal{H}\!\mathit{om}}\nolimits}
\def\SheafExt{\mathop{\mathcal{E}\!\mathit{xt}}\nolimits}
\def\Sch{\mathit{Sch}}
\def\Mor{\mathop{\mathrm{Mor}}\nolimits}
\def\Ob{\mathop{\mathrm{Ob}}\nolimits}
\def\Sh{\mathop{\mathit{Sh}}\nolimits}
\def\NL{\mathop{N\!L}\nolimits}
\def\CH{\mathop{\mathrm{CH}}\nolimits}
\def\proetale{{pro\text{-}\acute{e}tale}}
\def\etale{{\acute{e}tale}}
\def\QCoh{\mathit{QCoh}}
\def\Ker{\mathop{\mathrm{Ker}}}
\def\Im{\mathop{\mathrm{Im}}}
\def\Coker{\mathop{\mathrm{Coker}}}
\def\Coim{\mathop{\mathrm{Coim}}}

% Boxtimes
%
\DeclareMathSymbol{\boxtimes}{\mathbin}{AMSa}{"02}

%
% Macros for moduli stacks/spaces
%
\def\QCohstack{\mathcal{QC}\!\mathit{oh}}
\def\Cohstack{\mathcal{C}\!\mathit{oh}}
\def\Spacesstack{\mathcal{S}\!\mathit{paces}}
\def\Quotfunctor{\mathrm{Quot}}
\def\Hilbfunctor{\mathrm{Hilb}}
\def\Curvesstack{\mathcal{C}\!\mathit{urves}}
\def\Polarizedstack{\mathcal{P}\!\mathit{olarized}}
\def\Complexesstack{\mathcal{C}\!\mathit{omplexes}}
% \Pic is the operator that assigns to X its picard group, usage \Pic(X)
% \Picardstack_{X/B} denotes the Picard stack of X over B
% \Picardfunctor_{X/B} denotes the Picard functor of X over B
\def\Pic{\mathop{\mathrm{Pic}}\nolimits}
\def\Picardstack{\mathcal{P}\!\mathit{ic}}
\def\Picardfunctor{\mathrm{Pic}}
\def\Deformationcategory{\mathcal{D}\!\mathit{ef}}
\usepackage{cleveref}
\crefname{originallemma}{original lemma}{original lemmas}
\Crefname{originallemma}{Original lemma}{Original lemmas}
\makeatletter
\newlabel{lemma-cohomology-of-closed}{{09V3}{source}{}{sourcepointer-1}{}}
\newlabel{lemma-cohomology-of-closed@cref}{{[originallemma][0][]09V3}{[1][0][]source}{}{}{}}
\makeatother
\crefname{originallemma}{original lemma}{original lemmas}
\Crefname{originallemma}{Original lemma}{Original lemmas}
\makeatletter
\newlabel{lemma-proper-base-change}{{09V5}{source}{}{sourcepointer-2}{}}
\newlabel{lemma-proper-base-change@cref}{{[originallemma][0][]09V5}{[1][0][]source}{}{}{}}
\makeatother
\crefname{originaltheorem}{original theorem}{original theorems}
\Crefname{originaltheorem}{Original theorem}{Original theorems}
\makeatletter
\newlabel{theorem-proper-base-change}{{09V6}{source}{}{sourcepointer-3}{}}
\newlabel{theorem-proper-base-change@cref}{{[originaltheorem][0][]09V6}{[1][0][]source}{}{}{}}
\makeatother
\title{Proper base change for sheaves of sets on topological spaces}
\author{The Stacks Project Authors}
\date{stacks-project:cohomology.tex:local-snapshot-20261011; GNU FDL 1.2 or later}
\begin{document}
\maketitle
\noindent\textit{Editorial note.} All delivered mathematical text is the author's own, copied from the pinned Stacks Project chapter file cohomology.tex; no statement or proof step was rewritten, reordered or removed.\par\medskip
\noindent\textit{Editorial note.} The proof ends with the author's own 'Some details omitted.' after the reduction; the sentence is reproduced verbatim. The two reductions are named explicitly and the only input used is one whose first part is checked to work for sheaves of sets, so the omission is bounded.\par\medskip
\noindent\textit{Editorial note.} The counted spans carry 4 navigation stubs and no bibliography citation.\par\medskip
\noindent\textit{Editorial note.} The preamble is the chapter's retained minimal preamble (sources/stacks-cohomology/retained-preamble.tex lines 1-81), a verbatim excerpt of the native preamble.\par\medskip
\noindent\textit{Editorial note.} The label lemma-proper-base-change-sheaves-of-sets is defined only in cohomology.tex, so no cross-chapter restatement had to be weighed for this card.\par\medskip
\noindent\textit{Allowed standard prerequisites (editorial).}\begin{itemize}
\item Proper base change for complexes of abelian sheaves. Same chapter, Theorem theorem-proper-base-change and Lemma lemma-proper-base-change, whose proofs are followed at proof lines 3157-3161. Supplies the reduction to the stalk formula and the shape of the argument.
\item The zeroth degree case of the cohomology of a closed subset, valid for sheaves of sets. Same chapter, Lemma lemma-cohomology-of-closed, invoked at proof lines 3161-3163. Supplies the identification of the sections over the fibre.
\end{itemize}\medskip
\ifdefined\hypertarget\hypertarget{sourcepointer-1}{}\fi
\noindent\textit{Source reference.} Original lemma 09V3: sources/stacks-cohomology/pinned-original-source.tex, source lines 2848--2848. Pointer outside the counted span, invoked once for context; the load-bearing argument of the counted statement is the author's proof reproduced in full in the proof span.\par\medskip
\ifdefined\hypertarget\hypertarget{sourcepointer-2}{}\fi
\noindent\textit{Source reference.} Original lemma 09V5: sources/stacks-cohomology/pinned-original-source.tex, source lines 3059--3059. Pointer outside the counted span, invoked once for context; the load-bearing argument of the counted statement is the author's proof reproduced in full in the proof span.\par\medskip
\ifdefined\hypertarget\hypertarget{sourcepointer-3}{}\fi
\noindent\textit{Source reference.} Original theorem 09V6: sources/stacks-cohomology/pinned-original-source.tex, source lines 3101--3101. Pointer outside the counted span, invoked once for context; the load-bearing argument of the counted statement is the author's proof reproduced in full in the proof span.\par\medskip
\begin{lemma}[Proper base change for sheaves of sets]
\label{lemma-proper-base-change-sheaves-of-sets}
Consider a cartesian square of topological spaces
$$
\xymatrix{
X' \ar[d]_{f'} \ar[r]_-{g'} & X \ar[d]^f \\
Y' \ar[r]^g & Y
}
$$
Assume that $f$ is proper. Then
$g^{-1}f_*\mathcal{F} = f'_*(g')^{-1}\mathcal{F}$
for any sheaf of sets $\mathcal{F}$ on $X$.
\end{lemma}

\begin{proof}
We argue exactly as in the proof of Theorem \ref{theorem-proper-base-change}
and we find it suffices to show
$(f_*\mathcal{F})_y = \Gamma(X_y, \mathcal{F}|_{X_y})$.
Then we argue as in Lemma \ref{lemma-proper-base-change}
to reduce this to the $p = 0$ case of Lemma \ref{lemma-cohomology-of-closed}
for sheaves of sets. The first part of the proof of
Lemma \ref{lemma-cohomology-of-closed}
works for sheaves of sets and this finishes the proof.
Some details omitted.
\end{proof}

\end{document}
